\section*{Розділ V. Контроль та оскарження}
\addcontentsline{toc}{section}{Розділ V. Контроль та оскарження}

\subsection*{5.1. Право на оскарження}
\addcontentsline{toc}{subsection}{5.1. Право на оскарження}
    5.1.1. Право оскаржувати рішення, дії або бездіяльність ТВК мають:

        \begin{enumerate}[label=\alph*)]
            \item Кандидати на виборах, щодо яких прийнято оскаржуване рішення;
            \item Довірені особи кандидатів;
            \item Групи студентів у складі не менше 5 осіб, які мають обґрунтований інтерес у справі.
        \end{enumerate}

    5.1.2. Оскарженню підлягають:

        \begin{enumerate}[label=\alph*)]
            \item Рішення про відмову в реєстрації кандидата;
            \item Рішення про скасування реєстрації кандидата;
            \item Порушення процедури голосування;
            \item Порушення процедури підрахунку голосів;
            \item Рішення про визнання виборів недійсними або дійсними;
            \item Інші рішення ТВК, які порушують виборчі права студентів або суперечать цьому Регламенту.
        \end{enumerate}

    5.1.3. Не підлягають оскарженню організаційні рішення ТВК, що не впливають на результати виборів або виборчі права студентів.

\subsection*{5.2. Порядок подання скарг}
\addcontentsline{toc}{subsection}{5.2. Порядок подання скарг}
    5.2.1. Скарга на рішення, дії або бездіяльність ТВК подається у письмовій формі (в електронному вигляді з електронним підписом або у паперовому вигляді) протягом 2 (двох) робочих днів з моменту прийняття оскаржуваного рішення або виявлення порушення.

    5.2.2. Скарга повинна містити:

        \begin{enumerate}[label=\alph*)]
            \item Відомості про скаржника (П.І.Б., статус, контактні дані);
            \item Чітке формулювання оскаржуваного рішення, дії або бездіяльності;
            \item Обґрунтування незаконності або необґрунтованості оскаржуваного рішення з посиланням на норми цього Регламенту;
            \item Докази, що підтверджують доводи скарги;
            \item Вимоги скаржника (скасування рішення, проведення повторних виборів тощо);
            \item Дату подання та підпис скаржника.
        \end{enumerate}

    5.2.3. Скарга подається Голові ТВК, який протягом 4 годин з моменту отримання передає її для розгляду згідно з процедурою, визначеною цим Розділом.

    5.2.4. Скарги, подані з порушенням встановленого строку або без дотримання вимог до форми та змісту, залишаються без розгляду. Про це скаржник повідомляється протягом 6 годин.

\subsection*{5.3. Розгляд скарг Загальними зборами студентів}
\addcontentsline{toc}{subsection}{5.3. Розгляд скарг Загальними зборами студентів}
    5.3.1. Скарги на рішення ТВК розглядаються Загальними зборами студентів Університету, які скликаються спеціально для цієї мети.

    5.3.2. Загальні збори для розгляду скарги скликаються протягом 3 робочих днів з моменту надходження скарги. Повідомлення про скликання зборів поширюється через всі доступні офіційні канали комунікації не пізніше ніж за 24 години до початку зборів.

    5.3.3. Кворум для загальних зборів, що розглядають скарги, становить 100 студентів або 10\% від загальної кількості студентів Університету (береться менша цифра).

    5.3.4. У разі недосягнення кворуму протягом години після оголошеного часу початку зборів, збори переносяться на наступний день о тій самій годині з кворумом 50 студентів.

    5.3.5. Загальні збори розглядають скаргу за такою процедурою:

        \begin{enumerate}[label=\alph*)]
            \item Доповідь скаржника (до 10 хвилин);
            \item Пояснення представника ТВК (до 10 хвилин);
            \item Питання та обговорення (до 30 хвилин);
            \item Голосування щодо задоволення скарги.
        \end{enumerate}

\subsection*{5.4. Рішення за результатами розгляду скарг}
\addcontentsline{toc}{subsection}{5.4. Рішення за результатами розгляду скарг}
    5.4.1. За результатами розгляду скарги Загальні збори студентів можуть прийняти одне з наступних рішень:

        \begin{enumerate}[label=\alph*)]
            \item Залишити скаргу без задоволення;
            \item Задовольнити скаргу частково або повністю;
            \item Скасувати оскаржуване рішення ТВК;
            \item Зобов'язати ТВК прийняти нове рішення з урахуванням зауважень;
            \item Призначити повторні вибори (у випадках грубих порушень виборчого процесу).
        \end{enumerate}

    5.4.2. Рішення приймається простою більшістю голосів від присутніх на зборах студентів.

    5.4.3. Рішення Загальних зборів за скаргою є остаточним та обов'язковим для виконання ТВК.

    5.4.4. ТВК зобов'язана виконати рішення Загальних зборів протягом 24 годин з моменту його прийняття, якщо інший строк не встановлений у рішенні.

    5.4.5. Рішення Загальних зборів оформлюється протоколом, який підписується головуючим на зборах та секретарем, і оприлюднюється на офіційних інформаційних ресурсах протягом 6 годин.

\subsection*{5.5. Контроль за дотриманням виборчого процесу}
\addcontentsline{toc}{subsection}{5.5. Контроль за дотриманням виборчого процесу}
    5.5.1. Контроль за дотриманням процедур, встановлених цим Регламентом, здійснюють:

        \begin{enumerate}[label=\alph*)]
            \item Члени ТВК під час засідань комісії;
            \item Довірені особи кандидатів під час виборчого процесу;
            \item Спостерігачі від студентської спільноти;
            \item Члени дільничних виборчих комісій під час голосування та підрахунку голосів.
        \end{enumerate}

    5.5.2. Спостерігачі мають право:

        \begin{enumerate}[label=\alph*)]
            \item Бути присутніми на всіх відкритих засіданнях ТВК;
            \item Спостерігати за процесом голосування та підрахунку голосів;
            \item Робити фото- та відеофіксацію процедур (за згодою ТВК);
            \item Звертатися до ТВК з питаннями щодо дотримання процедур;
            \item Отримувати копії протоколів засідань та результатів голосування.
            \item Невідкладно здійснювати фото- та відеофіксацію подій, що на думку спостерігача містять ознаки порушення виборчого процесу, без попередньої згоди ТВК, за умови одночасного дотримання таких вимог:
                \begin{enumerate}[label=\roman*)]
                    \item не розкривати таємницю голосування (заборонено знімати заповнення бюлетенів, виборчі кабіни, екрани систем, списки виборців та інші носії персональних даних);
                    \item не перешкоджати роботі ТВК і дільничних комісій (у тому числі не блокувати проходи, місця роботи, огляд; не використовувати спалах та гучний звук; мінімізувати переміщення під час підрахунку);
                    \item негайно повідомити ТВК про факт такої фіксації та її підставу (короткий опис події);
                    \item при оприлюдненні або передачі матеріалів забезпечити знеособлення персональних даних (маскування облич, номерів документів, контактів тощо);
                    \item тривала зйомка процедур, пряма трансляція (стрім), а також зйомка робочих матеріалів, що містять персональні дані, здійснюються виключно за згодою ТВК.
                \end{enumerate}
        \end{enumerate}

    5.5.3. Спостерігачі не мають права втручатися в роботу ТВК або дільничних комісій, перешкоджати виборчому процесу, проводити агітацію.

    5.5.4. У разі виявлення порушень спостерігачі звертаються до ТВК з вимогою їх усунення та можуть ініціювати подання скарги через суб’єктів, визначених у п. 5.1.1, надаючи наявні матеріали фіксації.